% ---------------------------------------------------------------------------
% What is open, as physics.  Written 2026-10-24 by the Paper 1 stream.
%
% SCOPE, and it was wrong in the first draft.  This section carries open
% PHYSICS questions and what would settle each.  It does not carry register
% mechanics, numbering custody, or which stream owns what -- that is
% governance and belongs in the bundle's own .md files.  The first draft of
% this file imported all three and has been stripped.
%
% The register behind it is provenance/OPEN-QUESTIONS-v1.0.0.md, cited once
% below as a pointer and not summarised.
% ---------------------------------------------------------------------------
\section{What is open}
\label{sec:open-questions}

\emph{Each row is a physics question this release does not answer, with what
would settle it. The working history of each is in}
\texttt{provenance/OPEN-QUESTIONS-v1.0.0.md}.

\begin{center}\small
\begin{tabularx}{\textwidth}{@{}l >{\raggedright\arraybackslash}X >{\raggedright\arraybackslash}X@{}}
\toprule
 & \textbf{the question} & \textbf{what would settle it} \\
\midrule
\textbf{Q5} & Does the full $1+3$ treatment predict excess power at low
  $\ell$? & A rank-0/2 share that rises at low $\ell$ rather than with it.
  \textbf{F2}'s third panel measures it and finds the share flat at the
  quadrupole, which parks the question rather than answering it \\
\textbf{Q9} & The discretisations are audited; the \emph{physical} ranges are
  not & A stated validity range per approximation, with the measurement that
  bounds it \\
\textbf{Q10} & Does the second-order quadrupole ring for tensor modes? & The
  second-order source carried at rank two with a tensor initial condition ---
  recorded as a mechanism, not claimed \\
\textbf{Q11} & The acoustic source has the right phase and the spectrum does
  not & Downgraded from blocker to residual, \emph{not closed}: the oscillator
  tracks $\cos(kr_s)$ in exact phase and period, so the phase is lost between
  the source and the transfer \\
\textbf{Q12} & Two implementations of (102) agree at $k=1$ and diverge above
  $k=3$ & Which of the two reproduces the integral (102) is derived from ---
  \S\ref{sec:w-form} answers this for the bulk and not for the boundary \\
\textbf{Q13} & Which construction of $(\kappa'\tilde v_B)'$ is right? & The
  by-parts identity, term by term, against each construction \\
\textbf{Q14} & Should $C_0$ enter the line-of-sight integrand at all? &
  \S\ref{sec:ell-zero-row} settles the $\ell=0$ row; whether the resulting
  $C_0$ belongs in the bulk is the remaining half \\
\midrule
\textbf{Q15} & \textbf{What does omitting helium recombination cost?} The
  recombination history here integrates the hydrogen Peebles equation. Helium
  recombines earlier and shifts $x_e$, hence $\kappa'$, $\eta_*$, the
  visibility width and the damping scale. & $\eta_*$ and $k_D$ recomputed with
  helium, against the percent-level differences the peak comparison turns on.
  \emph{Not measured} \\
\textbf{Q16} & \textbf{Is the tight-coupling source frozen late enough?}
  $\delta T$ and $\tau_1$ are held at $\eta_{\rm couple}=6\eta_*$, which is
  correct physics --- after decoupling nothing drives the monopole --- on the
  condition that the visibility has closed by then. & $\mathcal V(\eta_{\rm
  couple})/\mathcal V(\eta_*)$ on every background. \emph{Nothing currently
  measures it}, and the recombination history differs between the three \\
\bottomrule
\end{tabularx}
\end{center}

\paragraph{Q16 cannot move a peak, and that is worth saying so it is not
over-read.} What is frozen is the tight-coupling pair, which after decoupling
has no scattering to drive it; the line-of-sight integral itself runs to
$\eta_0$ on a grid whose density is set by the Bessel kernel's own period in
$\eta$ and gated against aliasing --- finding \textbf{T-3}. A freeze that was
too early would cut source while the visibility was still live, which is a
boundary effect and not a projection one.

\paragraph{One open finding belongs beside these, because the plateau depends
on it.} \textbf{T-67} is open: $(a\mathcal B_1)'$ carries
$-a\mathcal H\,\delta T-3a\mathcal H\,\Phi_A$ beyond the terms quoted in
\S\ref{sec:w-form}, and whether they belong is unsettled. They are
$O(a\mathcal H)$ and the late integrated Sachs--Wolfe effect lives where
$k\sim a\mathcal H$, \textbf{so the reduced form is unavailable in exactly the
regime the low-$\ell$ plateau occupies.} Every statement of
$(a\mathcal B_1)'$ in this supplement is the expansion-terms-off form and says
so at its own site.

\paragraph{And the late ISW is a weight question before it is a physics one.}
The effect is already in the formalism: $\Phi'$ enters at rank one through
$(a\mathcal B_1)'$, hence through the bulk bracket. In any line-of-sight
assembly it rides on $W=e^{-\kappa}$ --- the photon accumulates $\int\Phi'$
along the whole path, weighted only by the probability it has not scattered
--- and \eqref{eq:byparts} agrees. \textbf{(102) carries that bracket at
$\mathcal V\mathcal D$, and $\mathcal V$ has closed by the time the late ISW
acts.} So discrepancy~(i) of \S\ref{sec:w-form} does not merely insert a
spurious $\kappa'$; it removes the late ISW. \emph{This is testable on the
background axis}: $\Phi$ is constant for Einstein--de~Sitter ($K=0$ dust,
Appendix~G of the antecedent), so EdS has no late ISW to lose, while
$\Lambda$CDM and \texttt{PLANCK\_MODEL} do. Correcting the weight should lift
the plateau on the latter two and leave EdS alone.

\textbf{That test has now run and the first half of it is confirmed.} The
late-time share of $C_\ell$ --- the fraction accumulated after the visibility
closes --- is \textbf{$0.0000$ in all thirty-six curves}: six multipoles,
three backgrounds, both frames. Einstein--de~Sitter is the control and
$0.0000$ is correct there; $\Lambda$CDM and \texttt{PLANCK\_MODEL} require
$\mathcal O(0.1)$ and return the same zero. The accumulation reproduces the
released quadrature at $1.15\times10^{-12}$, so it is that integral decomposed
and not a second definition of it. \textbf{A fault in the accumulation would
annihilate late deposition on every background including the one where zero is
correct, and a fault manufacturing it would show on EdS; neither signature is
present.} The paragraph above was written before the measurement and is
confirmed by it.

\textbf{What remains open is the second half}, and it is one run on an
existing control: with the bracket moved to $W$, the late share at $\ell=2$
must become non-zero on both $\Lambda$ backgrounds \emph{and must stay
exactly $0.0000$ on Einstein--de~Sitter}. \textbf{If it stays zero on the
$\Lambda$ backgrounds, discrepancy~(i) is not the mechanism and a second
effect shares its signature.} That is the discriminator, and it is what turns
the reading below from consistent-with into established.
